\documentclass[output=paper,%modfonts,nonflat,
colorlinks,citecolor=brown,
% hidelinks,
showindex
]{langscibook}
\ChapterDOI{10.5281/zenodo.18233968}
\author{Felix Anker\affiliation{University of Jena}}
\title{Negation in Tsova-Tush}  
\abstract{The present paper is a survey of negative constructions in the Nakh-Daghestanian language Tsova-Tush. Standard negation in main clauses, subordinate clauses and stative predications is expressed by the preverbal particle \textit{co}. This negator is also found in the formation of other negative constructions like negative pronouns, or in tag questions, as well as in negative replies. Another negation particle \textit{ma} is used with optatives and imperatives in order to form prohibitive constructions. Furthermore, there are several devices that can be used to derive negative items like privatives and to reinforce negation.

\keywords{negation, Tsova-Tush, negative pronouns, privative formation, negative lexicalization}
}


\IfFileExists{../localcommands.tex}{%hack to check whether this is being compiled as part of a collection or standalone
  % add all extra packages you need to load to this file

\usepackage{tabularx,multicol}
\usepackage{url}
\urlstyle{same}

\usepackage{listings}
\lstset{basicstyle=\ttfamily,tabsize=2,breaklines=true}

% \usepackage{langsci-basic}
\usepackage{langsci-optional}
\usepackage{langsci-lgr}
\usepackage{langsci-gb4e}

\usepackage[linguistics,edges]{forest}
\usepackage{subfigure}

\usepackage{arydshln}

\usepackage{pifont}
\newcommand{\cmark}{\ding{51}}

\usepackage{tabto}


\usepackage{multirow} %for Tsova-Tush, NGT, panara
\usepackage{booktabs} %for Tsova-Tush, panara
\usepackage{graphicx} %for Tsova-Tush
\usepackage{longtable} %for Tsova-Tush, NGT
\usepackage{vcell} %for Tsova-Tush

% \usepackage{pdflscape} %kambaata -- has been added to enable landscape formatting of table 3
\usepackage{xcolor} %kambaata,ngarinjin,nivkh -- has been added to enable gray shading in tables
\definecolor{light-gray}{gray}{0.8}%ngarinjin
\definecolor{lightest-gray}{gray}{0.95}%ngarinjin

%\usepackage{subcaption} %for NGT, but does not work with subfigure package from above
\usepackage{slgloss} %for NGT
% \newfontfamily{\handshape}{handshape.TTF} %for NGT %HC2022-11-04: it does not work and creates an error "Package fontspec Error: The font "handshape" cannot be found", eventhough I have uploaded the file handshape.TTF, so I outcommented it for the time being. Could you solve this Sebastian? thanks :)
%\newcommand\customfont[1]{{\usefont{T1}{handshape}{m}{n} #1 }} %for NGT
\newfontfamily\ToneFont{CharisSIL-Regular.ttf}
\AdditionalFontImprint{Charis SIL}
\usepackage[shortlabels]{enumitem} %NGT 


\usepackage{leipzig} %for the glosses in panara
\usepackage{glossaries} %panara
% \usepackage{lscape} %panara
% \input{libertine-preamble} %panara
% \setmathfont{Libertinus Math} %panara
%   \usepackage{enumerate} %panara - it clashes with enumitem which is really needed for NGT, and it works in panara also with enumitem, i.e. we don't need enumerate.

% \setmainfont[Ligatures=TeX,Numbers=OldStyle]{Linux Libertine O} %panara
% \setsansfont[Ligatures=TeX,Numbers=OldStyle]{Linux Biolinum O} %panara

\usepackage{hyperref}%kalamang,yukana
% \usepackage{tipx}%kalamang,ungarinjin, yukana

% \usepackage{wrapfig}%ungarinjin
% \usepackage{tabulary}%ungarinjin
%\usepackage{subfig}%ungarinjin %deleted bc subfigure already above
\PassOptionsToPackage{colorinlistoftodos,disable}{todonotes}
\usepackage{todonotes}
\usepackage{hyphenat}
\usepackage{langsci-branding}


  %\newcommand*{\orcid}{}

\makeatletter
\let\thetitle\@title
\let\theauthor\@author
\makeatother

\newcommand{\togglepaper}[1][0]{
   \bibliography{../localbibliography}
   \papernote{\scriptsize\normalfont
     \theauthor.
     \titleTemp.
     To appear in:
     Matti Miestamo \& Ljuba Veselinova (ed.).
     Title of book.
     Berlin: Language Science Press. [preliminary page numbering]
   }
   \pagenumbering{roman}
   \setcounter{chapter}{#1}
   \addtocounter{chapter}{-1}
}


\renewbibmacro*{index:name}[5]{% %Tsova-Tush
  \usebibmacro{index:entry}{#1} %Tsova-Tush
    {\iffieldundef{usera}{}{\thefield{usera}\actualoperator}\mkbibindexname{#2}{#3}{#4}{#5}}} %for Tsova-Tush

    
\newcommand{\gl}[1]{\textsc{#1}}

\newcommand{\ipa}[1]{{\textit{#1}}} %geshiza 
\newcommand{\ipaEX}[1]{{\textit{#1}}} %geshiza


%\newcommand{\hspaceThis}[1]{\hphantom{#1}} %this was in localcommands in the kambaata chapter

\newcommand{\ANA}{\textsc{ana}} %gyeli
\newcommand{\ATTR}{\textsc{attr}} %gyeli
\newcommand{\CONJ}{\textsc{conj}} %gyeli
\newcommand{\ID}{\textsc{id}} %gyeli
\newcommand{\IDEO}{\textsc{ideo}} %gyeli
\newcommand{\INCH}{\textsc{inch}} %gyeli
\newcommand{\LINK}{\textsc{link}} %gyeli
\newcommand{\NP}{NP} %gyeli
\newcommand{\NCA}{\textsc{nca}} %gyeli
\newcommand{\NPI}{\textsc{npi}} %gyeli
\newcommand{\PCF}{\textsc{pcf}} %gyeli
\newcommand{\PN}{\textsc{pn}} %gyeli
\newcommand{\PREP}{\textsc{prep}} %gyeli
\newcommand{\PRIOR}{\textsc{prior}} %gyeli
\newcommand{\PRO}{\textsc{pro}} %gyeli
\newcommand{\PROSP}{\textsc{prosp}} %gyeli
\newcommand{\R}{\textsc{r}} %gyeli
\newcommand{\RETRO}{\textsc{retro}} %gyeli
\newcommand{\SEQU}{\textsc{sequ}} %gyeli
\newcommand{\STAMP}{\textsc{stamp}} %gyeli
\newcommand{\SUB}{\textsc{sub}} %gyeli
\newcommand{\TM}{TM} %gyeli
\newcommand{\V}{\textsc{v}} %gyeli


\newcommand{\hspaceThis}[1]{\hphantom{#1}} %for Kambaata chapter, helps for \db (which allows alignment of word and gloss and not without taking into account the "s)


%for Panara:
\renewbibmacro*{index:name}[5]{%
  \usebibmacro{index:entry}{#1}
    {\iffieldundef{usera}{}{\thefield{usera}\actualoperator}\mkbibindexname{#2}{#3}{#4}{#5}}}

 %panara
\newleipzig{ade}{ade}{adessive} %panara
\newleipzig{dir}{dir}{directional} %panara
\newleipzig{emph}{emph}{emphasis} %panara
\newleipzig{fnl}{fnl}{final} %panara
\newleipzig{nml}{nml}{nominal form}%panara
\newleipzig{ine}{ine}{inessive}  %panara 
\newleipzig{ints}{ints}{intensifier} %panara
\newleipzig{iter}{iter}{iterative} %panara
\newleipzig{plac}{plac}{pluractional} %panara


% \makeglossaries %panara %HC not sure what to do with these following lines, but it does not work yet. I have added the abbreviations in the same table is in other chapters, so this glossary function is not needed anymore.
% \newglossarystyle{myglosses}{%
%   \renewenvironment{theglossary}%
% {\begin{multicols}{2}\raggedright}
% {\end{multicols}}
%     \renewcommand*{\glossaryheader}{}
%     \renewcommand*{\glsgroupheading}[1]{}
%     \renewcommand*{\glsgroupskip}{}
%     \renewcommand*{\glsclearpage}{} 
% % set how each entry should appear:
%     \renewcommand*{\glossentry}[2]{
%     \noindent\makebox[7em][l]{\glstarget{##1}{\textsc{\glossentryname{##1}}}}
%         \glossentrydesc{##1}
%     }
%     \renewcommand*{\subglossentry}[3]{%
%         \glossentry{##2}{##3}
%     }
% }

\newcommand{\glopa}{\textsc{dem}} %kalamang
\newcommand{\glme}{\textsc{top}} %kalamang
\newcommand{\glte}{\textsc{nfin}} %kalamang
\newcommand{\glse}{\textsc{iam}} %kalamang
\newcommand{\glkin}{\textsc{vol}} %kalamang
\newcommand{\glteba}{\textsc{prog}} %kalamang
\newcommand{\gli}{\textsc{plnk}} %kalamang
\newcommand{\glet}{\textsc{irr}} %kalamang
\newcommand{\glta}{\textit{ta}} %kalamang
\newcommand{\glnn}{\textit{n}} %kalamang

\newcommand{\tabitem}{~~\llap{\textbullet}~~}%kogi

%For Matti and Ljuba's comments on the last version of the chapters:

\newcommand{\matti}[2][color=green!40,size=\small]{\todo[#1]{MM: #2}}
\newcommand{\ljuba}[2][color=blue!40,size=\small]{\todo[#1]{LV: #2}}
\newcommand{\mm}[2][color=green!40,size=\small]{\todo[#1]{MM: #2}}
\renewcommand{\lv}[2][color=blue!40,size=\small]{\todo[#1]{LV: #2}}
\newcommand{\hc}[2][size=\small]{\todo[#1]{HC: #2}}
\newcommand{\gloss}[2][color=red!30,size=\tiny]{\todo[#1]{Glossing: #2}}
\newcommand{\glosst}[2][color=red!30,size=\tiny,inline]{\todo[#1]{Glossing: #2}} %t for (in the) text / inline
\newcommand{\aut}[2][color=yellow!50,size=\tiny]{\todo[#1]{Author: #2}}%for author

\renewcommand{\keywords}[1]{%
  \par\noindent\textbf{Keywords: #1.}%
}


\renewcommand{\REF}[2][]{(\ref{#2}#1)}
\renewcommand{\sectref}[1]{§\ref{#1}}
\newcommand{\Abf}{{Ā}\kern-2pt\raisebox{3.1pt}{́}\kern2pt}
\newcommand{\abf}{{ā}\kern-1pt\raisebox{1.1pt}{́}\kern1pt}
\newcommand{\A}{{Ā}\kern-1.5pt\raisebox{3.1pt}{́}\kern1.5pt}
\newcommand{\U}{{ʊ}\kern-1.5pt\raisebox{-.4pt}{̀}\kern1.5pt}

\newcommand{\fdot}{{f\kern-.75pt\raisebox{-2.5pt}{̣}\kern.75pt}}

\newcommand{\handshape}[2][1.5]{\includegraphics[height=#1ex]{fonts/#2.png}}
 
  \togglepaper[23]
}{}

\begin{document}

\maketitle

\section{The language}
\label{tsova:introduction}
Tsova-Tush (ISO: \href{https://iso639-3.sil.org/code/bbl}{bbl}, Glottocode: \href{https://glottolog.org/resource/languoid/id/bats1242}{bats1242}), also known as Batsbi, belongs to the Nakh branch of the Nakh-Daghestanian (or Northeast Caucasian) language family, together with Chechen and Ingush. It is spoken by fewer than 500 people in one village in the Eastern part of the Republic of Georgia. The majority of speakers are older than 50 and children no longer learn the language. The Tsova-Tush people came from the mountains to the lowlands in the 19th century and since then have been heavily influenced by Georgian, not only culturally but also linguistically, which is reflected in the vast amount of loanwords but also in some grammatical constructions, see also \citet{WichersSchreur2025}. 

Tsova-Tush is a morphologically ergative language,  with some fluid intransitive verbs \citep{Holisky1987}. Beside a number of grammatical cases, there are a number of spatial cases. All cases, beside the absolutive which is zero marked, need a so-called oblique stem marker to which they are attached. Another salient feature of Tsova-Tush is gender, a feature that is found in most Nakh-Daghestanian languages. There is  a total of eight genders, and gender agreement is mainly triggered by the absolutive arguments, i.e. S and P arguments. Gender agreement is indexed by prefixes on many verbs, some adjectives, adverbs and one numeral and postposition.

%\footnote{Alternatively, the last three groups are analyzed as mixed genders, the singular belonging to one group, plural to another (see \citealt{HoliskyGagua1994})

Due to the influence of Georgian, Tsova-Tush additionally developed person agreement which is not found in many other Nakh-Daghestanian languages. First and second person S and A arguments trigger person agreement which is marked by a suffix that is undoubtably of pronominal origin. 

The data presented here come from different sources. By far the biggest corpus is the Endangered Caucasian Languages in Georgia  (ECLinG) corpus,\footnote{The ECLinG corpus is available in the MPI archive which can be accessed via \url{https://archive.mpi.nl/tla/islandora/object/tla\%3A1839_00_0000_0000_0008_24AD_F}} where most of the data stem from. More data is found in several Tsova-Tush  dictionaries (\citealt{Kadagidze1984}; \citealt{BertlaniEtAl2013}; \citealt{BertlaniEtAl2019}) and other existing literature on the language (e.g. \citealt{HoliskyGagua1994}, a short grammar sketch; \citealt{Schiefner1856}, the first description of the language). Additionally, I used data from my own fieldwork.

The focus of this paper is the expression of negation in Tsova-Tush and its structure follows the questionnaire provided by the editors \parencitetv{chapters/questionnaire}: \sectref{tsova:clausalnegation} deals with clausal negation, i.e. negation in declarative and non-declarative clauses, negation in stative predications and negation in non-main clauses. In \sectref{tsova:nonclausal} non-clausal negation is illustrated, that is, negative replies, negative indefinites and negative derivational devices. Other aspects of negation such as the scope of negation, reinforcing negation and coordination of negated clauses are treated in \sectref{tsova:otheraspects}. Finally \sectref{tsova:summary} will sum up the most important features of negation in Tsova-Tush.

\section{Clausal negation} 
\label{tsova:clausalnegation}

\subsection{Standard negation}
\label{tsova:standardnegation}

Standard negation in Tsova-Tush is quite regular and achieved by means of the negator \textit{co}. This particle always appears preverbally, as is expected for an SOV language, and there is no structural asymmetry (cf. \citealt{Miestamo2005}) between positive and negative clauses, i.e. they only differ in absence vs. presence of the negation marker as can be seen in examples (\ref{tsova:cartridge}) and (\ref{tsova:horseshoe}) which both contain the verb \textit{d-aɬ-iⁿ} `\textsc{iii}-give-\textsc{aor}'.\footnote{Roman numerals in the gloss are used to indicate gender.} In the first clause it is affirmative while in the second one it is used with the negator \textit{co}.

\ea	\label{tsova:cartridge}
\gll 	cħa     pxi p'at'run d-aɬ-iⁿ son=e o tapar-v daħ-v-ax-en-es.\\  
	one five cartridge\textsc{(iii).abs} \textsc{iii}-give-\textsc{aor} \textsc{1sg.dat=coord} this gun-\textsc{ins} \textsc{pvb-i}-go-\textsc{aor-1sg.erg}\\
\glt `He gave me some five cartridge cases and I went there with that gun.' (ECLinG, bav22\_11, 00:01:52)
\z

\ea	\label{tsova:horseshoe}
\gll 	erek'le mepe-s oquj-n nal-i co d-aɬ-iⁿ.\\  
	Erekle king-\textsc{erg}  \textsc{3sg}.\textsc{obl}-\textsc{dat} horseshoe\textsc{(iii).abs}-\textsc{pl} \textsc{neg} \textsc{iii}-give-\textsc{aor}\\
\glt `King Erekle didn't give him horseshoes.' (ECLinG, bav17\_06, 00:00:28)
\z

Tsova-Tush has several preverbs to express spatial and aspectual relations. In complex verbs that consist of a preverb and the lexical verb, the negation particle follows the preverb and precedes the verb, (\ref{tsova:goinside}).
This not only shows that the preverbs are not fully attached to the verb, but also that the negator \textit{co} could be treated as a prefix rather than a particle. There are, however, some other instances, where \textit{co} is not fully integrated into the verb, e.g. in polar interrogative clauses, the interrogative enclitic \textit{=i} attaches to the item in preverbal position, which is –  in negative polar interrogatives – \textit{co}. Therefore, I will treat it as a particle, as in \REF{tsova:goinside}.


\ea	\label{tsova:goinside}
\gll 	atx	uq	deni	ču	co	ix-atx.\\  
	\textsc{1pl.excl.erg} that.\textsc{obl} day.\textsc{in} \textsc{pvb}.in \textsc{neg} go-\textsc{1pl.excl.erg}\\
\glt `We don’t go inside that day.' (ECLinG, bav21\_05, 00:01:16)
\z

Tsova-Tush has periphrastic verb forms consisting of a participle and the auxiliary \textit{ daⁿ} `be'. The negator \textit{co} tends to precede the entire construction, i.e. it occurs either before the participle, as in (\ref{tsova:heardsong}), or it occurs before the auxiliary, as in (\ref{tsova:flower}). Its occurence before the auxiliary is much less common. 


\ea	\label{tsova:heardsong}
\gll 	ištruⁿ	vum	co	xac’-en	d-a	soⁿ  ǯer.\\  
	such thing\textsc{(iii).abs} \textsc{neg} hear-\textsc{ptcp} \textsc{iii}-be.\textsc{prs} \textsc{1sg.dat} yet\\
\glt `I haven’t heard such a thing yet.' (ECLinG, bav19\_06, 00:01:03)
\z

\ea	\label{tsova:flower}
\gll 	co j-a-ħ daħ-q'anjajln vʕalaʔ=a uišne bubuk'-saⁿ lamzur j-a-ħ.\\  
	\textsc{neg} \textsc{ii}-be.\textsc{prs}-\textsc{2sg.abs} \textsc{pvb}-become\_old.\textsc{ptcp} at\_all=\textsc{coord} similar  flower-like beautiful  \textsc{ii}-be.\textsc{prs}-\textsc{2sg.abs}\\
\glt `You have not become old at all, you are as beautiful as a flower.' (ECLinG, bav14\_07, 00:00:13)
\z

In complex verbs that consist of a root and a transitivizing or intransitivizing derivational suffix, \textit{co} cannot intervene between the root and the suffix (cf.  \citet{Harris2009}), e.g.  (\ref{tsova:lambs}).

There is no inherent `not yet' form in the sense of \citet[54–55]{Comrie1985} in Tsova-Tush. Expression of such constructions is achieved by the Georgian particle \textit{ǯer} `yet', as in (\ref{tsova:yet}).

\ea	\label{tsova:yet}
\gll 	seⁿ badr-i-v ǯer co aɬin d-a so-go.\\
	\textsc{1sg.gen} child-\textsc{pl-erg} 	yet \textsc{neg} say.\textsc{ptcp} \textsc{iii}-be.\textsc{prs} \textsc{1sg-all.in}\\
\glt `My children haven't said that to me yet.' (ECLinG, bav13\_17, 00:05:16)
\z

A suffix that is only found in negative contexts is \textit{-gĕ} which is attached to verbs to express `no longer, not anymore' and is used together with the negation particle \textit{co}. In present and aorist forms it follows the tense suffixes,  as in (\ref{tsova:anymore}), while the past tense suffix and personal markers follow it, as in (\ref{tsova:nolongerlaugh}--\ref{tsova:nolonger}).

\ea	\label{tsova:anymore}
\gll 	soⁿ	ħo	co	j-ec’e-gĕ.\\  
	\textsc{1sg.dat} \textsc{2sg.abs} \textsc{neg} \textsc{ii}-love.\textsc{prs}-anymore\\
\glt `I don't love you (female) anymore.' (fieldnotes)
\z

\ea	\label{tsova:nolongerlaugh}
\gll 	co j-il-ge-r.\\  
	\textsc{neg} \textsc{ii}-laugh-anymore-\textsc{pst}\\
\glt `She was no longer laughing.' \citep[187]{HoliskyGagua1994} 
\z

\ea	\label{tsova:nolonger}
\gll 	co aɬ-in-ge-s.\\  
	\textsc{neg} say-\textsc{aor}-anymore-\textsc{1sg.erg}\\
\glt `I no longer said it.' \citep[187]{HoliskyGagua1994} 
\z

From time to time the suffix \textit{-gĕ} is marked twice on verbs. According to my consultant this double marking  is used to draw attention to the emotional bond the speaker has toward the event, as in (\ref{tsova:anymoreanymore}).

\ea	\label{tsova:anymoreanymore}
\gll 	ese alni duq daxn co d-a-ge-gĕ.\\  
	here Alvani\textsc{.in} many cattle\textsc{(iii).abs} \textsc{neg}  \textsc{iii}.be.\textsc{prs}-anymore-anymore\\
\glt `Here in Alvani, there are not many cattle anymore.' (ECLinG, bav07\_13, 00:04:36)
\z

\subsection{Negation in non-declaratives}
\label{tsova:nondec}

Non-declarative clauses, i.e. interrogative and imperative clauses show some deviation from the standard negation pattern illustrated in \sectref{tsova:standardnegation}. 

In content questions the negation particle \textit{co} behaves as in standard negation, it simply precedes the verb as illustrated in (\ref{tsova:why}) with the interrogative pronoun \textit{vuⁿ} `why'. 

\ea	\label{tsova:why}
\gll 	vux	aɬ-in,	vuⁿ	co	d-axk’-eⁿ.\\  
	what say-\textsc{aor} why \textsc{neg} \textsc{iii}-come-\textsc{aor}\\
\glt `What did they say, why didn’t they come back?' (ECLinG, bav19\_16, 00:01:47)
\z

Polar questions, however, are marked by the interrogative enclitic =\textit{i} which is attached to the item that precedes the verb. In case of negation this is \textit{co}, which –  through a regular phonological process –  becomes \textit{cu=j}, as in (\ref{tsova:deserter}). If, however, the focus of the polar question is on another constituent, it will take the interrogative enclitic instead, as can be seen in (\ref{tsova:CHILD}), where the focus of the question is on \textit{bader} `child'.

\ea	\label{tsova:deserter}
\gll 	dezert’	cu=j	v-a?\\  
	deserter\textsc{(i).abs} \textsc{neg=q} \textsc{i}-be.\textsc{prs}\\
\glt `Isn't he a deserter?' (ECLinG, bav24\_09, 00:00:51)
\z

\ea	\label{tsova:CHILD}
\gll 	bader=i co d-a?\\  
	child\textsc{(iii).abs}=\textsc{q} \textsc{neg} \textsc{iii}-be.\textsc{prs}\\
\glt `Isn't it a \textsc{child}?' (fieldnotes)
\z

Prohibitive constructions are formed by using the negation particle \textit{ma}, which precedes the imperative. Just like the standard negator \textit{co} it occurs in preverbal position, as in (\ref{tsova:pimple}), and after preverbs, as in (\ref{tsova:dontforget}). Positive and negative imperatives only differ in presence vs. absence of the prohibitive particle, i.e. they belong to group II of \citegen{AuweraEtAl2013} classification of prohibitives: “The prohibitive uses the verbal construction of the second singular imperative and a sentential negative strategy
not found in (indicative) declaratives.”

\ea	\label{tsova:pimple}
\gll 	seⁿ ag bʕarče is levan-ate-r [k'ak'al duq ma j-aq'o-t coħek'a kort-i-x muc'uk' d-ejɬ šun=ajnŏ].\\  
	\textsc{1sg.gen} grandmother always this say-\textsc{hab}-\textsc{pst} nut(\textsc{iv}).\textsc{abs} many \textsc{proh} \textsc{iv}-eat-\textsc{imp.pl} otherwise head-\textsc{obl-ccs} pimple(\textsc{iii}).\textsc{abs} \textsc{iii}-appear \textsc{2pl.dat=quot}\\
\glt `My grandmother always used to say: “Don't eat too many nuts, otherwise pimples will break out on your head."' (ECLinG, bav04\_01, 00:08:33)
\z

\ea	\label{tsova:dontforget}
\gll 	daħ	ma	j-ic-j-o-l	jaš-o!\\  
	\textsc{pvb} \textsc{proh} \textsc{ii}-forget-\textsc{ii-intrz-pol} sister-\textsc{voc}  
\\
\glt `Sister, don't forget!'  \citep[146]{Mikeladze2011}\footnote{There is no vocative in Tsova-Tush; the suffix has been borrowed from Georgian}
\z

In negative optatives, which express wishes, curses, and the like and are formed by adding the politeness suffix \textit{-l} to the imperative suffix, the same particle \textit{ma} is used, as in (\ref{tsova:seesun}).

\ea	\label{tsova:seesun}
\gll 	matx	ma	gib-a-l	oquj-n!\\  
	sun \textsc{proh} see-\textsc{imp-pol} \textsc{3sg.obl-dat}\\
\glt `May s/he not see the sun!' \citep[182]{HoliskyGagua1994} 
\z

\subsection{Negation in stative predications}
\label{tsova:stative}

Negation in stative predications does not differ from standard negation. All stative predications, i.e. equation, proper inclusion, attribution, locative predication, existential predication and possessive predication, are headed by the copula. If negated, the standard clausal negator \textit{co} precedes the copula, as in (\ref{tsova:ball}–\ref{tsova:donkey}), and, again, there is no structural asymmetry between negative and positive stative predications. Unlike periphrastic verbal construcions, where  \textit{co} precedes the entire construction, in stative predications, the negator precedes the copula, not the predicate.

Although \citet[63]{Schiefner1856} mentions the omission of the copula-final vowel when used with \textit{co}, i.e. \textit{co d-a} `it is not (\textsc{neg} \textsc{iii}-be.\textsc{prs})' becomes \textit{cod}, I could not find examples of this form in my data. According to my speaker, this form is not used anymore even though she notes that her grandmother, born in the late 19th century, used this form. Later generations, however, only used it when mockingly imitating the speech of older people.\footnote{Another similar contraction that is still in use is \textit{comad} from \textit{com} `nothing' (see \ref{tsova:pronouns}) and \textit{d-a} `\textsc{iii}-be.\textsc{prs}'}

\ea	\label{tsova:ball}
Equation\\
\gll 	o	juq’	co	j-a.\\  
	this ball(\textsc{iv}).\textsc{abs} \textsc{neg} \textsc{iv}-be.\textsc{prs}\\
\glt `This is not a ball.' (ECLinG, bav23\_08, 00:03:35)
\z

\ea	\label{tsova:somemen}
Proper Inclusion\\
\gll 	meⁿ st'ak' st'ak' co v-a\\  
	some man\textsc{(i).abs} man\textsc{(i).abs} \textsc{neg} \textsc{i}-be\textsc{.prs}\\
\glt `Some men are not men.' \citep[147]{BertlaniEtAl2013} 
\z

\ea	\label{tsova:bighouse}
Attribution\\
\gll 	e c'a d-aqːoⁿ co d-a.\\  
	that house\textsc{(iii).abs} \textsc{iii}-big \textsc{neg} \textsc{iii}-be\textsc{.prs}\\
\glt `This house is not big.' (fieldnotes)
\z

\ea	\label{tsova:spider}
Locative Predication\\
\gll 	keǯbera co d-a č’er-e-x.\\  
	spider(\textsc{iii}).\textsc{abs} \textsc{neg} \textsc{iii}-be.\textsc{prs} wall-\textsc{obl-ccs}\\
\glt `The spider is not on the wall.' (fieldnotes)
\z

\ea	\label{tsova:nowork}
Existential Predication\\
\gll 	normalurat	ise	samušao	co	j-a.\\  
	normally here work(\textsc{iv}).\textsc{abs} \textsc{neg} \textsc{iv}-be.\textsc{prs}\\
\glt `Normally, there is no work here.' (ECLinG, bav21\_10, 00:02:38)
\z

\ea	\label{tsova:donkey}
Possessive Predication\\
\gll 	vir	co	d-a	txo-go.\\  
	donkey(\textsc{iii}).\textsc{abs} \textsc{neg} \textsc{iii}-be.\textsc{prs} \textsc{1pl.excl-all.in}\\
\glt `We don't have a donkey.' (ECLinG, bav07\_01, 00:11:11)
\z

\subsection{Negation in non-main clauses}
\label{tsova:nonmain}

Tsova-Tush non-main clauses are headed by a number of non-finite verb forms like converbs to express temporal and other adverbial relations, participles, or verbal nouns (masdars). In negative subordinate clauses, these are preceded by the negation particle \textit{co}. 

Examples (\ref{tsova:sad}) and (\ref{tsova:lambs}) illustrate negative adverbial clauses, marked by the simultaneous converb and the conditional converbs, respectively, with the preposed negator \textit{co}.

\ea	\label{tsova:sad}
\gll 	[čuħ	co	j-a-š]	daɣonbadjailnŏ	j-a-sŏ.\\ 
	home \textsc{neg} \textsc{ii}-be.\textsc{prs}-\textsc{sim.cvb} sad.\textsc{ii} \textsc{ii}-be.\textsc{prs}-\textsc{1sg.abs}\\
\glt `When I'm not at home, I'm sad.' \citep[203]{HoliskyGagua1994} 
\z

\ea	\label{tsova:lambs}
\gll 	[xi co maɬ-d-i-čeħ-a], [xi-n mak co ħerc'-d-i-čeħ-a] čuxui vomaʔ daħ-daxašle-bala ħo-go.\\  
	water \textsc{neg} drink-\textsc{iii-tr}-\textsc{cond.cvb}-\textsc{2sg.erg} spring-\textsc{dat} on \textsc{neg} turn-\textsc{iii-tr}-\textsc{cond.cvb}-\textsc{2sg.erg} lamb\textsc{(iii).pl} all \textsc{pvb}-die\_out-\textsc{intrz} \textsc{2sg-all.in}\\
\glt `If you don't give the lambs water, if you don't let them walk around a spring, all of them will die on you.' (ECLinG, bav03\_02, 00:06:54)
\z

Negative adverbial clauses that are formed by means of the verbal noun behave just as the negative converbal adverbial clauses, i.e. the negator precedes the verbal noun as exemplified in (\ref{tsova:horse}) where the causal adverbial clause is headed by the verbal noun in the contact case which is preceded by the negator \textit{co}.

\ea	\label{tsova:horse}
\gll 	[doⁿ	co	xiɬ-r-ex]	kuitad	v-eʔ-n-as.\\  
	horse \textsc{neg} become-\textsc{msd}-\textsc{ccs} on\_foot \textsc{i}-come-\textsc{aor}-\textsc{1sg.erg}\\
\glt `Because of not having a horse, I came on foot.' \citep[204]{HoliskyGagua1994} 
\z

One of the strategies to form relative clauses in Tsova-Tush, as is typical for Nakh-Daghestanian languages, is by using participles. In (\ref{tsova:care}), the head of the relative clause is a participle which, in negative contexts, is preceded by the negation particle. 


\ea	\label{tsova:care}
\gll 	[že-x	daħ	co	v-ac’ujn]	st’ak’.\\
	sheep-\textsc{ccs} \textsc{pvb} \textsc{neg} \textsc{i}-care\_for.\textsc{pst.ptcp} man\textsc{(i).abs}\\
\glt `A man who didn’t care for sheep.' (ECLinG, bav24\_08, 00:01:13)
\z

While Chechen, one of the two sister languages of Tsova-Tush, has negative converbs that are used, for instance, in resultative constructions (cf. \citealt[146]{Molochieva2010}), there are no negative converbs in Tsova-Tush.

\subsection{Negative lexicalizations}
\label{tsova:neglex}

There is one negative lexicalization in Tsova-Tush, \textit{vuntxeʔ} `not know', as in (\ref{tsova:drankthat}). Note that the use of  \textit{vuntxeʔ} as a complement taking item is rare and it mainly functions as a parenthetical. A more frequent way of expressing `not know' is by means of the negator \textit{co} and the verb \textit{xeʔaⁿ} `know.\textsc{inf}'. 

\ea	\label{tsova:drankthat}
\gll 	oqus o xi ħal-maɬ-en, vuntxeʔ oqujn vux d-a.\\  
	\textsc{3sg.erg} that	water	\textsc{pvb}-drink-\textsc{aor}	not.know	\textsc{3sg.dat}		what \textsc{iii}-be.\textsc{prs}\\
\glt `She drank that water not knowing what it was.'  (ECLinG, bav18\_02, 00:01:49)
\z

\section{Non-clausal negation}
\label{tsova:nonclausal}

\subsection{Negative replies}
\label{tsova:reply}

Negative replies to positive polar questions are expressed by the negation particle \textit{co}. It can either function on its own as a prosentence as in example \REF{tsova:no} or be accompanied by additional information. 

\ea	\label{tsova:no}
    \begin{xlist}[B: ]
    \exi{Q: }
        \gll osiħ=i d-ex-ne-š? \\  
        	there=\textsc{q} \textsc{iii}-live-\textsc{aor}-\textsc{2pl} \\
        \glt `Did you live there?' 
    \exi{A: }
        \gll co.\\  
        	\textsc{neg}\\
        \glt `No.' (ECLinG, bav22\_16, 00:02:04) 
    \end{xlist}
\z 

The same particle \textit{co}  is used in order to provide a negative reply to a negative polar question, as in (\ref{tsova:love}). Repetition of the whole clause is also possible in negative replies although it is much less common. An example for complete repetition is  (\ref{tsova:dontwork}) where the negated clause is used without an additional preceding negative particle \textit{co}. The only difference is the shift in person marking from second person to first person.

\ea	\label{tsova:love}
\gll 	dec’ar	vux	d-a	cu=j	xeʔ	ħoⁿ?	co=en	aɬ-iⁿ.\\  
	love(\textsc{iii}).\textsc{abs} what \textsc{iii}-be.\textsc{prs} \textsc{neg=q} know.\textsc{prs} \textsc{2sg.dat} \textsc{neg=quot} say-\textsc{aor}\\ 
\glt `Don't you know what love is? “No", she said.' (ECLinG, bav19\_02, 00:00:52)
\z

\ea	\label{tsova:dontwork}
    \begin{xlist}[Q: ]
    \exi{Q: }
        \gll 	co	mušebado? \\  
        	\textsc{neg} work.\textsc{prs.2sg.abs} \\	
        \glt `Don’t you work?'
    \exi{A: }
        \gll co	mušebado-s.\\  
        	\textsc{neg} work.\textsc{prs-1sg.abs}\\	
        \glt `I don’t work.' (ECLinG, bav08\_04, 00:00:40)
    \end{xlist}
\z 

If there is an affirmative answer to a negative polar question there is a polarity-reversing particle \textit{moħec}, whose primary meaning is `of course', but it can be used to contradict the negative question, similar to special particles that are found in other languages which fulfill the same function, e.g. German \textit{doch} or Hungarian \textit{de} (for an extensive overview on polarity reversing particles cf. \citealt{Holmberg2016}, and \citealt{Moser2019}). In (\ref{tsova:harmonica}), \textit{moħec} is used in this function expressing the speaker's familiarity with the harmonica in contrast to the propositional content of the question.

\ea	\label{tsova:harmonica}
    \begin{xlist}[Q: ]
    \exi{Q: }
        \gll 	lungɣar nat'uiⁿ garmoⁿ cu=j xac'en j-a ħoⁿ? \\  
        	 Longianti Nato.\textsc{gen} harmonica(\textsc{iv}).\textsc{abs}. \textsc{neg=q} hear.\textsc{ptcp} \textsc{iv}-be.\textsc{prs} \textsc{2sg.dat} \\	
        \glt `Didn't you listen to Logianti Nato's harmonica?'
    \exi{A: }
        \gll moħec, xac'en j-a soⁿ.\\  
        	of\_course hear.\textsc{ptcp} \textsc{iv}-be.\textsc{prs} \textsc{1sg.dat}\\	
        \glt `Yes, of course I listened to it.' (ECLinG, bav24\_12, 00:03:29)
    \end{xlist}
\z 

The impact of the Georgian language can be seen in negative replies as well since the use of the Georgian negation particle \textit{ara} is found occasionally, as in (\ref{tsova:sell}).

\ea	\label{tsova:sell}
\begin{xlist}[Q: ]
    \exi{Q: }
        \gll 	 	daħ-b-oxk’-aⁿ	co	leʔ-er	ħoⁿ?	 \\  
        	 \textsc{pvb-v}-sell-\textsc{inf} \textsc{neg} wish-\textsc{pst} \textsc{2sg.dat}\\	
        \glt `Didn’t you want to sell it?'
    \exi{A: }
        \gll ara,	co	b-ujxk’-n-as.\\  
        	no \textsc{neg} \textsc{v}-sell-\textsc{aor-1sg.erg}\\	
        \glt `No, I didn’t sell it.' (ECLinG, bav05\_10, 00:01:13)
    \end{xlist}
\z

\subsection{Negative indefinites and quantifiers}
\label{tsova:pronouns}

Negative indefinites in Tsova-Tush are interrogative-based (cf. \citealt{Haspelmath1997,Haspelmath2013-wals46}) and are formed by combining the standard negation marker \textit{co} with interrogative pronouns. The only exception to this is the negative indefinite \textit{com} `nothing' which is a contraction of the standard negator  \textit{co} and the indefinite pronoun \textit{vum} `something'. A partial paradigm is given in \tabref{tsova:tab:pron} which shows the formation of negative indefinites with `who' and `what'. As can be seen in \tabref{tsova:tab:pron}, the interrogative pronouns are suppletive, i.e. absolutive pronouns are based on the stems \textit{meⁿ} and \textit{vux}, while the oblique pronoun stems are \textit{ħan-} and \textit{st'en-}.\footnote{This form is also found in answers to `Thank you', i.e. \textit{cost'eni} `You're welcome!'}

\begin{table}
\caption{Tsova-Tush interrogative and negative absolutive and oblique pronouns}
\label{tsova:tab:pron}
\begin{tabular}{@{}llllll@{}}
\lsptoprule
             & \multicolumn{2}{c}{`who'}       &  & \multicolumn{2}{c}{`what'}          \\
             \cmidrule{2-3}\cmidrule{5-6}
             & Interrogative & Negative        &  & Interrogative   & Negative          \\
             \midrule
\textsc{abs} & \textit{meⁿ}  & \textit{comena} &  & \textit{vux}    & \textit{com}      \\
\textsc{obl} & \textit{ħan-} & \textit{coħan-} &  & \textit{st'en-} & \textit{cost'en-} \\
\textsc{dat} &	 \textit{ħann} & \textit{coħanna} &  & \textit{st'enn} & \textit{cost'enna} \\
\lspbottomrule
\end{tabular}
\end{table}

Further negative indefinites are `nowhere' and `never' which are formed in a similar way. The combination of the standard negator \textit{co} and the interrogative \textit{mičeħ} `where' yields \textit{comičħe} `nowhere'. This negative indefinite can be furthermore modified by spatial case markers. In (\ref{tsova:fromnowhere}), the ablative marker is suffixed to the negative indefinite. Similarly, the negating particle \textit{co} can combine with the interrogative \textit{macaⁿ} `when' which results in the negative indefinite  \textit{comacni}, as in (\ref{tsova:cantfind}), although the expression of  `never' is most of the time achieved by means of simple standard negation with the particle \textit{co}. 

\ea	\label{tsova:fromnowhere}
\gll 	iq dač'irba sašvelav comič-rena j-a so-go.\\  
	this.\textsc{obl} trouble.\textsc{in} help\textsc{(iv).abs} nowhere-\textsc{abl} \textsc{iv}-be.\textsc{prs} \textsc{1sg-all.in}\\
\glt `In this trouble, help is coming from nowhere.' \citep[272]{BertlaniEtAl2013} 
\z

\ea	\label{tsova:cantfind}
\gll 	comacni lax-mak' soⁿ ħo.\\  
	never find-\textsc{pot} \textsc{1sg.dat} \textsc{2sg.abs}\\
\glt `I can never find you.' (fieldnotes)
\z

According to \citet[202ff]{Mikeladze2011} there are three types of clauses with negative indefinites, one of them being the original Nakh one while the others are borrowings from Georgian and mirror the historical development of Georgian negative indefinite clauses. The three patterns are as follows:

\begin{enumerate}
    \item the use of an indefinite pronoun in a negative clause which seems to be the original Nakh pattern, as in (\ref{tsova:dogbark1}), and (\ref{tsova:work}) for comparison to Ingush 
    \item the use of a negative indefinite pronoun in an affirmative clause, as in (\ref{tsova:dogbark2}--\ref{tsova:allowed})
    \item the use of a negative pronoun in a negative clause, i.e. double negation, as in (\ref{tsova:dogbark3}--\ref{tsova:booksaw}) 
\end{enumerate}

In the first pattern, an indefinite pronoun is used in a negated clause. This strategy is the original strategy in Nakh languages but seems to be in the course of being replaced by the Georgian patterns. Compare examples (\ref{tsova:dogbark1}) from Tsova-Tush and (\ref{tsova:work}) from Ingush, its sister language. There are no dedicated negative indefinite pronouns in Ingush but if used in negative clauses, the generic-noun-based indefinites \textit{sag} `person' and \textit{hama} `thing' function as negative pronouns (cf. \citealt[696]{Nichols2011}). 

\ea	\label{tsova:dogbark1}
\gll 	ħam-en	co	xajc’-n-or	txe	pħarči-ⁿ	axa-r.\\  
	someone.\textsc{obl-dat} \textsc{neg} hear-\textsc{aor-rep} \textsc{1pl.gen} dog.\textsc{pl-gen} bark-\textsc{msd}\\
\glt `Nobody heard our dogs’ barking.' \citep[198]{Mikeladze2011} 
\z

\ea	\label{tsova:work} Ingush (East Caucasian, ISO: \href{https://iso639-3.sil.org/code/inh}{inh}, Glottocode: \href{https://glottolog.org/resource/languoid/id/ingu1240}{ingu1240})\\
\gll 	taxan balxa sag hwa-vendzar.\\  
	today work.\textsc{adv} person  \textsc{dx-v}.come.\textsc{neg.wit.pst}\\
\glt `Nobody came to work today.' \citep[696]{Nichols2011} 
\z

The second pattern, the use of a negative indefinite pronoun in an affirmative clause, seems to be a Georgian calque and presents a structurally intermediate stage between the lack of negative indefinite pronouns and negative concord. This strategy is also the most common in \citeauthor{Schiefner1856}'s early grammar of Tsova-Tush from 1856, although the third pattern is found occasionally. Examples (\ref{tsova:dogbark2}) and  (\ref{tsova:allowed}) illustrate the use of negative indefinite pronouns in affirmative clauses. Just as in Georgian, the position of the negative pronoun in this construction is restricted to the preverbal position. In the third pattern, negation is additionally marked by preverbal \textit{co}, allowing the negative pronoun to have a less restricted position (cf. for Georgian \citealt{HollowayKing1996}), as in (\ref{tsova:booksaw}).


\ea	\label{tsova:dogbark2}
\gll 	coħan-na xajc’-n-or	txe	pħarči-ⁿ	axa-r.\\  
	nobody.\textsc{obl-dat} hear-\textsc{aor-rep} \textsc{1pl.gen} dog.\textsc{pl-gen} bark-\textsc{msd}\\
\glt `Nobody heard our dogs’ barking.' \citep[202]{Mikeladze2011} 
\z

\ea	\label{tsova:allowed}
\gll 	isev comena d-ec'e-r d-aʔ-it-aⁿ oqus.\\  
	here nobody.\textsc{abs} \textsc{iii}-must-\textsc{pst} \textsc{iii}-come-\textsc{caus}-\textsc{inf} \textsc{3sg.erg}\\
\glt `She wasn't allowed to let anybody enter there.' (ECLinG, bav18\_04, 00:10:41)
\z

The use of negative pronouns in negative clauses is a rather recent innovation, probably a calque from Georgian. Constructions with this double negation do not differ from the previous two strategies in meaning. Example (\ref{tsova:dogbark3}) can be used to express the same proposition as examples (\ref{tsova:dogbark1}) and (\ref{tsova:dogbark2}). As mentioned above, negative pronouns or quantifiers do not need to be adjacent to the verb or the negator in this construction but can be followed by other constituents, as in (\ref{tsova:booksaw}). 

\ea	\label{tsova:dogbark3}
\gll 	coħa-na	co	xajc’-n-or	txe	pħarči-ⁿ	axa-r.\\  
	nobody.\textsc{obl-dat} \textsc{neg} hear-\textsc{aor-pst} \textsc{1pl.gen} dog.\textsc{pl-gen} bark-\textsc{msd}\\
\glt `Nobody heard our dogs’ barking.' \citep[202]{Mikeladze2011} 
\z

\ea	\label{tsova:booksaw}
\gll 	comičħe ħeⁿ žagnŏ co d-ag-ir soⁿ\\  
	nowhere \textsc{2sg.gen} book\textsc{(iii).abs} \textsc{neg} \textsc{iii}-see-\textsc{pst} \textsc{1sg.dat}\\
\glt `I haven't seen your book anywhere.' (fieldnotes)
\z

Negative pronouns in prohibitive constructions and optatives are used in the same manner although the standard negator \textit{co} is replaced by the prohibitive particle \textit{ma}, as in (\ref{tsova:nobodycome}).

\ea	\label{tsova:nobodycome}
\gll 	mamena	ma	d-o-la-l	esev.\\  
	nobody.\textsc{abs} \textsc{proh} \textsc{iii}-come-\textsc{imp-pol} here\\
\glt `Nobody come here!' (fieldnotes)
\z

In polar questions which contain negative indefinite pronouns, the position of the interrogative marker \textit{=i} varies. With \textit{com} `nothing' the interrogative marker is attached to the negative pronoun, i.e. \textit{com=i}, as in (\ref{tsova:knowanything}). On the other hand, in polar questions that involve \textit{comena}  `nobody', the interrogative marker is attached to the negator, i.e. it is placed between the negator with which the negative indefinite is formed and the interrogative pronoun, \textit{cu-j-mena} `\textsc{neg-q}-who', as in (\ref{tsova:speaknobody}).

\ea	\label{tsova:knowanything}
\gll 	com=i	xeʔ	soⁿ.\\  
	nothing\textsc{=q} know.\textsc{prs} \textsc{1sg.dat}\\
\glt `Don't you know anything?' (ECLinG, bav16\_02, 00:00:40)
\z

\ea	\label{tsova:speaknobody}
\gll 	psare cu-j-mena lav-iⁿ šu-gŏ?\\  
	yesterday \textsc{neg}-\textsc{q}-who talk-\textsc{aor} \textsc{2pl-all}\\
\glt `Did nobody talk to you yesterday?' (fieldnotes)
\z

Further negative indefinites are formed  with \textit{com} `nothing'  and are used mainly to reinforce negation (see \sectref{tsova:reinforcing}), \textit{com-ač'} `nothing at all (nothing.\textsc{abs-ints})', \textit{com-c} `not at all (nothing.\textsc{abs}-indeed)', \textit{cħa-ʔ-com} `nothing at all (one-again-nothing)', \textit{com-gĕ} `nothing anymore (nothing-anymore)', \textit{com-c-ge-ʔ} `nothing anymore (nothing-{indeed}-anymore-again)', \textit{cħa-ʔ-co} `not even one (one-again-\textsc{neg})' and \textit{cħa-comena} `nobody at all (one-nobody.\textsc{abs})'. 

\subsection{Negative derivation and case marking}
\label{tsova:derivation}

There are several derivational processes in Tsova-Tush to form privatives.
The primary strategy to form privatives is by means of the suffix \textit{-c'iⁿ} which can be attached to nouns, as in (\ref{tsova:master}), and verbal stems, as in (\ref{tsova:uneducated}). 

\ea	\label{tsova:master}
\gll 	txeⁿ c'a dadi-c'iⁿ d-is-en=e išt'aʔ d-a-r-atx.\\  
	\textsc{1pl.excl.gen} house\textsc{(iii).abs} master-\textsc{priv} \textsc{iii}-remain-\textsc{aor}=\textsc{coord} such \textsc{iii}-be-\textsc{pst}-\textsc{1pl.excl.abs}\\
\glt `Our house remained without a master and so we lived.' (ECLinG, bav19\_13, 00:01:58)
\z

\ea	\label{tsova:uneducated}
\gll 	ʡamda-c’iⁿ	j-a-r=e	p’ap’ra	gor-e-ⁿ	joħ	j-a-r.\\  
	learn-\textsc{priv} \textsc{ii}-be-\textsc{pst}=\textsc{coord} Papashvili family-\textsc{obl-gen} girl\textsc{(ii).abs} \textsc{ii}-be-\textsc{pst}\\
\glt `She was uneducated and she was a girl from the Papashvili family.' (ECLinG, bav22\_08, 00:00:18)
\z

Another possibility to express privatives is direct borrowing of Georgian privatives that are formed by the circumfix \textit{u- -o} with loss of the word-final vowel. While this strategy is restricted to Georgian loanwords, the Tsova-Tush privative suffix \textit{-c'iⁿ} can be applied not only to indigenous Tsova-Tush words but to Georgian loanwords as well. In the first example below, Tsova-Tush has borrowed the privative form of the Georgian word with loss of the suffix. In the second example, only the noun has been borrowed from Georgian, and the Tsova-Tush privative suffix is added.

\begin{enumerate}
\item Georgian \textit{ut'vino} > Tsova-Tush \textit{ut'viⁿ}  `brainless' 
\item Georgian \textit{uk'almo} vs. Tsova-Tush \textit{k'almec'iⁿ}  `without a pen' 
\end{enumerate}

With some adjectives it is possible to derive antonyms by prefixing the standard negation marker \textit{co} to it, e.g. \textit{coɣazen} `not good, bad' from \textit{ɣazen} `good'. 

\section{Other aspects of negation}
\label{tsova:otheraspects}
\subsection{The scope of negation}
\label{tsova:scope}

In focused negative constructions, the scope of negation can be narrowed to the NP in focus. The NP in the scope of negation is followed by the negator \textit{co} which – in standard negation –  precedes the verb and can be additionally marked by \textit{ma}  `however', as in (\ref{tsova:mybrother}).

\ea	\label{tsova:mybrother}
\gll 	mit'ŏ seⁿ vaš  v-a=e p'et'ŏ ma co.\\  
	Mito \textsc{1sg.gen} brother\textsc{(i).abs} \textsc{i}-be.\textsc{prs=coord} Peto however \textsc{neg}\\
\glt `Mito is my brother, not Peto.' \citep[210]{HoliskyGagua1994} 
\z

\subsection{Negative polarity items}

Beside the negative pronouns that are only found in negative constructions there are two more items that are only used in negative contexts. The first one is \textit{vʕalaʔ} `at all' which is used to reinforce negation (\sectref{tsova:reinforcing}) The second one is the postposition \textit{bedeⁿ} `but, besides, except for' which is treated in \sectref{tsova:nonneguse}. Although the resulting clause is does not express a negative meaning, it still requires the clause to be marked by the negation particle \textit{co}. 

\subsection{Marking of NPs in the scope of negation}

There are no effects on the marking of NPs in negative constructions. 

\subsection{Reinforcing negation}
\label{tsova:reinforcing}
\largerpage

There are several items that are used to reinforce negation. A frequent one is the negative polarity item \textit{vʕalaʔ} which would best be translated as `at all'. It is only found in negative contexts and can be used with the standard negation marker \textit{co},  as in (\ref{tsova:flower2}), as well as with the prohibitive marker \textit{ma}, as in (\ref{tsova:noise}), and negative pronouns, as in (\ref{tsova:hair}). Its position seems not to be fixed.

\ea	\label{tsova:flower2}
\gll 	co j-a-ħ daħ-q'anjajln vʕalaʔ=a uišne bubuk'-saⁿ lamzur j-a-ħ...\\  
	\textsc{neg} \textsc{ii}-be.\textsc{prs}-\textsc{2sg.abs} \textsc{pvb}-become\_old.\textsc{ptcp} at.all=\textsc{coord} similar  flower-like beautiful  \textsc{ii}-be.\textsc{prs}-\textsc{2sg.abs}\\
\glt `You have not become old at all, you are as beautiful as a flower.' (ECLinG, bav14\_07, 00:00:13)
\z

\ea	\label{tsova:noise}
\gll 	vʕalaʔ tataⁿ ma j-aqo-t, aɬ-iⁿ o badri-g, at'in ʡedaxk'-a-t, aɬ-iⁿ...\\  
	at.all noise\textsc{(iv).abs} \textsc{proh} \textsc{iv}-spread-\textsc{imp.pl}, say-\textsc{aor} this children-\textsc{all}, quiet sit.\textsc{pl-imp-imp.pl} say-\textsc{aor}\\
\glt `“Don't make any noise," he said to the children, “sit there silently ..."' (ECLinG, bav18\_04, 00:15:57)
\z

\ea	\label{tsova:hair}
\gll 	j-axːeⁿ beǯ j-a o Maqvle-go-ħ=e, lamzur zorejš=e, šer nana-s vʕalaʔ comičħe ix-it-or.  \\
	\textsc{iv}-long hair\textsc{(iv).abs} \textsc{iv}-be.\textsc{prs} this Maqvala-\textsc{all-in=coord} beautiful very\textsc{=coord} \textsc{refl} mother-\textsc{erg} at.all nowhere.\textsc{in} go-\textsc{caus-pst}\\
\glt `That Maqvala has long hair, very beautiful and her mother did not allow her to go anywhere at all.' (ECLinG, bav18\_02, 00:00:06)
\z

Another strategy to reinforce negation is the use of the negative reinforcing item \textit{comcgeʔ} (\textit{com-c-ge-ʔ} `nothing-{indeed}-anymore-again')  that is used to express `nothing anymore'.  Similarly, \textit{comgĕ} (\textit{com-gĕ} `nothing-anymore') can be used to fulfill the same function. 
Further items found with \textit{com} `nothing'  to reinforce negation are \textit{comač'} `nothing at all' which uses the intensifying suffix \textit{-ač'}, as in (\ref{tsova:relax}) and \textit{com-c} `not at all (nothing-{indeed})'.

\ea	\label{tsova:relax}
\gll 	ma kotːl-a-t, tiv-a-t=ajnŏ xajc'nŏ, baq'eʔ comač' tag-din-or den-cħaʔ. \\
	\textsc{proh} sorrowful.\textsc{intrz}-\textsc{imp-imp.pl} relax-\textsc{imp-imp.pl=quot} understand.\textsc{aor} really nothing\_at\_all do-\textsc{tr}-\textsc{pst} day-one\\
\glt `“Don't worry, relax!", they understood and they did nothing at all for the whole day.'  (fieldnotes)
\z

Other negative pronouns that are based on the numeral \textit{cħa} `one' in order to express `at all' are e.g.   \textit{cħa-ʔ-co} `not even one (one-again-\textsc{neg})', \textit{cħa-ʔ-com} `nothing at all (one-again-nothing)', as in (\ref{tsova:inlaw}) and \textit{cħa-ʔ-comena}  `nobody at all (one-again-nobody)', as in (\ref{tsova:servants}).

\ea	\label{tsova:inlaw}
\gll 	inc marnan-i-v cħaʔcom d-ec' aɬ-aⁿ.\\  
	now mother\_in\_law-\textsc{pl-erg} nothing\_at\_all \textsc{iii}-must.\textsc{prs} say-\textsc{inf}\\
\glt `Now, the mothers-in-law mustn't say anything at all.' (ECLinG, bav16\_06, 00:05:20)
\z

\ea	\label{tsova:servants}
\gll 	qeⁿ aɬ-iⁿ o k'nat-i-v, msaxur-i-v me cħaʔcomena co d-a=en aɬ-iⁿ.\\
	then say-\textsc{aor} that boy-\textsc{pl-erg} servant-\textsc{pl-erg} \textsc{sub} nobody\_at\_all \textsc{neg} \textsc{iii}-be.\textsc{prs=quot} say-\textsc{aor}\\
\glt `Then those boys, the servants, said ``Nobody is here at all''.' (ECLinG, bav18\_04, 00:14:17)
\z

The same strategy to reinforce negation is found in prohibitives, where \mbox{\textit{cħaʔcom}} is replaced by its prohibitive counterpart \textit{cħaʔmam}, as in (\ref{tsova:donttake}).

\ea	\label{tsova:donttake}
\gll 	cħaʔmam d-ik'-a-l=ajnŏ, d-ex-iⁿ vaš-a-s nan-a-x.\\
	nothing\_at\_all \textsc{iii}-take-\textsc{imp-pol=quot} \textsc{iii}-request-\textsc{aor} brother-\textsc{obl-erg} mother-\textsc{obl-ccs}\\
\glt `“Don't take anything at all'', the brother asked his mother.' (fieldnotes) 
\z
    
Finally, there is a suffix, \textit{-k'aʔ} `a little', which is usually used with adjectives or adverbs (e.g. \textit{kast'en-k'aʔ}  `a little quickly (quickly-\textsc{dim})'). If used with the negator \textit{co} it expresses a rather negative attitude of the speaker towards the proposition although not as drastic as standard negation, as in (\ref{tsova:meat}).
    

\ea	\label{tsova:meat}
\gll 	o ħal-qexk'-ien co-k'aʔ j-ec' soⁿ.\\
	this \textsc{pvb}-boil-\textsc{ptcp} \textsc{neg-dim} \textsc{iv}-love.\textsc{prs} \textsc{1sg.dat}\\
\glt `I don't like it (the meat) too much when it is boiled up.' (ECLinG, bav08\_12, 00:08:20)
\z 
   

\subsection{Negation, coordination  and complex clauses}
\label{tsova:coordination}
There is no special negative coordinating conjunction for coordinating affirmative and negative clauses, or two negative clauses. This can be achieved by juxtaposition, by using the coordinator \textit{je} `and' or the coordinating enclitic \textit{=e} `and' respectively.

Another strategy that is found with coordinated negative clauses, as well as with negated coordinated constituents, is coordination by means of the particle \textit{le} `or' which is used to form negative correlative conjunctions. If used to coordinate negative clauses, the particle is placed clause initially and each clause is either negated independently or, if they share the same verb, only one clause receives a negator. 

In coordinated prohibitive clauses the same particle \textit{le} `or' is used but together with the prohibitive particle \textit{ma}. 

When used with negated coordinated constituents, \textit{le} precedes each constituent, as in (\ref{tsova:rotten}), while the negator is found in its standard position, i.e. preceding the verb. 

\ea	\label{tsova:rotten}
\gll 	txen pħeħ daħ co b-axk'l-ar xolme le qor, le sxal, le cħaʔcom.\\  
	our village.\textsc{in} \textsc{pvb} \textsc{neg} \textsc{v}-rot-\textsc{pst} use\_to or apple or pear or nothing\_at\_all\\
\glt `In our village neither the apples, nor the pears, nor anything else became rotten.' (ECLinG, bav04\_01, 00:17:57)
\z

Tsova-Tush has a negative conjunction, \textit{coħek'aʔ} `lest, otherwise', which is made up of the negator, the conditional \textit{-ħĕ}, and \textit{-k'aʔ} `a little', as in (\ref{tsova:pimple}) and (\ref{tsova:otherwise}).

\ea	\label{tsova:otherwise}
\gll 	so-g v-ek-iⁿ me ħal ma qitː  coħek'aʔ p'irdap'ir qos-o-s ħoⁿ aɬ-iⁿ.\\  
	\textsc{1sg-all} \textsc{i}-shout-\textsc{aor} \textsc{sub} \textsc{pvb} \textsc{proh} get\_up lest straight fire-\textsc{fut}-\textsc{1sg.erg} \textsc{2sg.dat} say-\textsc{aor}\\
\glt `He shouted at me: ``Don't stand up, otherwise I will fire at you."' (ECLinG, bav06\_07, 00:02:37)
\z

An example for contrastive negation can be found in \sectref{tsova:scope}.

Tsova-Tush has one modal verb to express possibilities, permission and the like, \textit{mak'aⁿ} `can, be able', which can appear in two constructions.

\begin{itemize}
	\item \textit{mak’aⁿ}  can be attached directly to the verbal stem, as in (\ref{tsova:beer})
	\item \textit{mak’aⁿ } is used finitely together with an infinitive, as in (\ref{tsova:cantsing}--\ref{tsova:gothere})
\end{itemize}

In negation with the first pattern the negator \textit{co} precedes the verbal stem to which the potential marker is attached, as in (\ref{tsova:cantgo2}--\ref{tsova:beer}). Tense and person markers follow the potential suffix.

\ea	\label{tsova:cantgo2}
\gll 	co	v-ax-mak’	soⁿ.	\\  
	\textsc{neg} \textsc{i}-go-\textsc{pot.prs} \textsc{1sg.dat}\\
\glt `I can’t go.' (fieldnotes)
\z

\ea	\label{tsova:beer}
\gll 	co	xaʔ-mak’	soⁿ	[ʕuⁿ	meɬŏ	vaser-v	jeg].\\  
	\textsc{neg} know-\textsc{pot.prs} \textsc{1sg.dat} why drink.\textsc{prs} men-\textsc{erg} beer\\
\glt `I can’t understand why men drink beer.' (fieldnotes)
\z

In the second pattern, i.e. with \textit{mak’aⁿ } used as a finite verb, the negator can either precede the potential verb, i.e. it negates the matrix clause, as in (\ref{tsova:cantsing}), or it precedes the infinitive inside the subordinate clause, as in (\ref{tsova:gothere}).

\ea	\label{tsova:cantsing}
\gll 	co	mak’	ħoⁿ	[inc	moqb-aⁿ].\\  
	\textsc{neg} can.\textsc{prs} \textsc{2sg.dat} now sing-\textsc{inf}\\
\glt `You can’t sing now.' (fieldnotes)
\z

\ea	\label{tsova:gothere}
\gll 	[ču	co 	v-ax-aⁿ]	mak’	soⁿ me [bin  miɣeba-d-o-lo-s].\\  
	\textsc{pvb} \textsc{neg} \textsc{i}-go-\textsc{inf} can.\textsc{prs} \textsc{1sg.dat} \textsc{sub} flat(\textsc{iii}).\textsc{abs} receive-\textsc{iii}-\textsc{tr}-\textsc{subj}-\textsc{1sg.erg}\\
\glt `I can’t go there to take the flat.' (ECLinG, bav25\_15, 00:03:53)
\z

In negative constructions, however, the first pattern is preferred. The second pattern is found rarely overall but seems to be preferred for affirmative constructions. The distribution of the two patterns in negative and affirmative contexts is illustrated in \tabref{tsova:tab:mak}.

\begin{table}
\caption{Distribution of \textit{mak’aⁿ } in negative and positive clauses}
\label{tsova:tab:mak}
\begin{tabular}{llll} 
\lsptoprule
                                    & NEG   & POS  & n  \\ \midrule
V-\textit{mak'aⁿ}                   & 91    & 5    & 96 \\
\textit{mak'aⁿ} + V-\textsc{inf}    & 6     & 21   & 27   \\ 
\lspbottomrule
\end{tabular}
\end{table}

\subsection{Miscellaneous aspects of negation}
\subsubsection{Negative transport}

There are at least two instances of negative transport, i.e.  the negator of the matrix clause is interpreted as negating the subordinate predicate. Those two items that allow negative transport are the verb \textit{motːaⁿ} `think' and \textit{tabit} which can best be translated as `it seems (to me)', as in (\ref{tsova:notthink}--\ref{tsova:tabitnot}). 

\ea	\label{tsova:notthink}
\gll 	co motː soⁿ [me o v-aɣ-ŏ].\\  
	\textsc{neg} think\textsc{.prs}  \textsc{1sg.dat} \textsc{sub}  \textsc{3sg}  \textsc{i}-come-\textsc{prs}\\
\glt `I don't think that he is coming.' (fieldnotes)
\z

\ea	\label{tsova:thinknot}
\gll 	motː soⁿ [me o co v-aɣ-ŏ].\\  
	think\textsc{.prs} \textsc{1sg.dat} \textsc{sub}	\textsc{3sg} \textsc{neg} \textsc{i}-come-\textsc{prs}\\
\glt `I think that he is not coming.' (fieldnotes)
\z

\ea	\label{tsova:tabit}
\gll 	co tabit [v-aɣ-ŏ o].\\  
	\textsc{neg} it\_seems \textsc{i}-come-\textsc{prs} \textsc{3sg}\\
\glt `It does not seem that he is coming.' (fieldnotes)
\z

\ea	\label{tsova:tabitnot}
\gll 	tabit [o co v-aɣ-ŏ]\\  
	it\_seems \textsc{3sg} \textsc{neg} \textsc{i}-come-\textsc{prs}\\
\glt `It seems that he is not coming.' (fieldnotes)
\z

It is matter of future research to establish whether other predicates allow negative transport.

\subsubsection{Non-negative uses of negatives and expletive negation}
\label{tsova:nonneguse}

Another  construction that involves the standard negation marker \textit{co} is the tag question. Tsova-Tush tag questions are formally the same as negation in negative polar interrogatives, namely, they are formed by combining the negative particle \textit{co} and the interrogative enclitic \textit{=i}. The resulting tag question marker \textit{cuj} follows the clause it refers to. In contrast to e.g. English tag questions, in Tsova-Tush there is no reversed polarity, i.e. the negative tag question is used with positive as well as negative clauses.  In (\ref{tsova:poor}), the clause that precedes the tag question marker \textit{cuj} is negated, while in (\ref{tsova:respect}) the main clause is positive. In affirmative clauses there is another possibility by using the particle \textit{haʔ=i} `yes=\textsc{q}', although the former strategy seems to be preferred with both polarities.

\ea	\label{tsova:poor}
\gll 	ħaʔ,	magram	oquj-go	co	j-a-r	gač’irveb,	cu=j?\\  
	yes but \textsc{3sg-all.in} \textsc{neg} \textsc{iv}-be-\textsc{pst} poverty\textsc{(iv).abs} \textsc{neg=q}\\
\glt `Yes, but this one wasn’t poor, was he?' (ECLinG, bav25\_16, 00:01:45)
\z

\ea	\label{tsova:respect}
\gll 	j-aqːoⁿ	p’ativ	j-ebl-or	marnan-e-n	mardad-e-n,	cu=j?\\  
	\textsc{iv}-big respect\textsc{(iv).abs} \textsc{iv}-put-\textsc{pst} mother\_in\_law-\textsc{obl-dat} father\_in\_law-\textsc{obl-dat} \textsc{neg=q}\\
\glt `They respected their mother-in-law and father-in-law, didn’t they?' (ECLinG, bav15\_06, 00:05:07)
\z

The postposition \textit{bedeⁿ} `but, besides, except for'  is another expression that is only found in formally negative contexts that actually encode a positive sense. In (\ref{tsova:octopus}), the postposition governs its noun in the case that is required by the verb, here the contact case. While the proposition itself is not negative, the postposition can only be used together with the negator \textit{co}.

\ea	\label{tsova:octopus}
\gll 	rvapeqaʔo surt-a-x bedeⁿ co d-ag-in d-a soⁿ.\\  
	octopus\textsc{(iii).abs} picture-\textsc{obl-ccs} besides \textsc{neg} \textsc{iii}-see-\textsc{ptcp} \textsc{iii}-be.\textsc{prs} \textsc{1sg.dat}\\
\glt `I have only seen an octopus in a picture.' \citep[243]{BertlaniEtAl2013} 
\z

\subsubsection{Diachronic notes and observations}

Since there is not a lot of historical data other than \citet{Schiefner1856}, and today there is no new generation speaking the language, there is not a lot of material for diachronic observations. We can still offer some reasoning based on the available data.

The first change we find is the expansion of the contraction in negative copula constructions. The formerly acceptable, morphologically complex forms \textit{cod, cov, coj, ...} which consist of the negator and the gender marker with loss of word-final vowel of the copula are now only available as independent word forms, i.e. \textit{co da, co va, co ja, ...} (see \sectref{tsova:stative}). 

Another observation concerns negative pronouns. While the language started out with using indefinite pronouns  in negative clauses, due to Georgian influence Tsova-Tush adopted negative indefinite pronouns which were first used in affirmative clauses and are now subject to negative concord, i.e. with additional negation of the clause in addition to negative indefinite marking (see \sectref{tsova:pronouns}). 

\section{Summary}
\label{tsova:summary}
This paper gave an overview on negation in Tsova-Tush. \tabref{tsova:tab:sum} summarizes the most important negation strategies. 

% \begin{table}
% \resizebox{\textwidth}{!}{%
% \centering
% \caption{Summary of negation in Tsova-Tush}
% \label{tsova:tab:sum}
% \begin{tabular}{llll}
% \lsptoprule
% \textbf{Clausal negation} & & Chapter & Examples\\ \midrule
% \textit{co} & 
%     \begin{tabular}[t]{@{}l@{}}preverbal particle used to negate\\ 
%         \hspace{4mm} (1) main clauses\\ 
%         \hspace{4mm} (2) subordinate clauses\\ 
%         \hspace{4mm} (3) stative predications
%     \end{tabular} & 
%     \begin{tabular}[t]{@{}l@{}}\phantom{X}\\ 
%         \sectref{tsova:standardnegation}\\ 
%         \sectref{tsova:nonmain}\\ 
%         \sectref{tsova:stative}
%     \end{tabular} & 
%     \begin{tabular}[t]{@{}l@{}}\phantom{X}\\ 
%         (\ref{tsova:horseshoe}–\ref{tsova:heardsong})\\ 
%         (\ref{tsova:sad}–\ref{tsova:care})\\ 
%         (\ref{tsova:ball}–\ref{tsova:donkey})
%     \end{tabular} \\
% \textit{ma} & 
%     \begin{tabular}[t]{@{}l@{}}preverbal particle used to negate \\ 
%         \hspace{4mm} (1) imperatives\\ 
%         \hspace{4mm} (2) optatives
%     \end{tabular} & 
%     \begin{tabular}[t]{@{}l@{}}\phantom{X}\\ 
%         \sectref{tsova:nondec}\\ 
%         \sectref{tsova:nondec}
%     \end{tabular} & 
%     \begin{tabular}[t]{@{}l@{}}\phantom{X}\\ 
%         (\ref{tsova:pimple}--\ref{tsova:dontforget})\\ 
%         (\ref{tsova:seesun})
%     \end{tabular} \\
% \textit{cuj} & 
%     \begin{tabular}[t]{@{}l@{}}
%         preverbal particle \textit{co} + interrogative enclitic \textit{=i}\\ 
%         used to negate polar interrogatives
%     \end{tabular} & 
% \sectref{tsova:nondec} & (\ref{tsova:deserter}--\ref{tsova:CHILD}) \\
% \textit{-gĕ} & 
%     \begin{tabular}[t]{@{}l@{}}
%         verbal suffix only used in negative clauses,\\ 
%         expresses `no longer'
%     \end{tabular} & 
% \sectref{tsova:standardnegation} & (\ref{tsova:anymore}–\ref{tsova:anymoreanymore})\\
% \textit{ǯer} & particle to express `not yet' in negative clauses & \sectref{tsova:standardnegation} & (\ref{tsova:yet}) \\
% \textit{vuntxeʔ} & negative lexicalization `not know' & \sectref{tsova:neglex} & (\ref{tsova:drankthat}) \\ 
% \midrule \rule{0pt}{15pt}
% \textbf{Non-clausal negation} &&& \\ 
% \midrule
% \textit{co} & 
%     \begin{tabular}[t]{@{}l@{}}
%         negative particle used as default negative \\ 
%         reply to positive and negative questions
%     \end{tabular}
% & \sectref{tsova:reply} & (\ref{tsova:no}--\ref{tsova:love})\\
% \textit{moħec}& 
%     \begin{tabular}[t]{@{}l@{}}
%         polarity reversing particle meaning `of course',\\ 
%         used as affirmative reply to negative questions
%     \end{tabular} & 
% \sectref{tsova:reply} & (\ref{tsova:harmonica}) \\
%     \begin{tabular}[t]{@{}l@{}}
%         \textit{co/ma} \\ 
%         + interrogative pronoun
%     \end{tabular} & formation of negative pronouns & 
% \sectref{tsova:pronouns} & \tabref{tsova:tab:pron} \\
% \textit{-c'iⁿ}& 
%     \begin{tabular}[t]{@{}l@{}}
%         suffix used for the formation of privatives \\ 
%         from nouns and verbal stems
%     \end{tabular} & 
% \sectref{tsova:derivation} & (\ref{tsova:master}--\ref{tsova:uneducated}) \\
%     \begin{tabular}[t]{@{}l@{}}
%         \textit{cuj}
%     \end{tabular} & 
%     \begin{tabular}[t]{@{}l@{}}
%         negative particle \textit{co} + interrogative enclitic \textit{=i}, \\ 
%         used to mark tag questions
%     \end{tabular} & 
% \sectref{tsova:nonneguse} & (\ref{tsova:poor}–\ref{tsova:respect}) \\ 
% \midrule \rule{0pt}{15pt}
% \textbf{Other aspects of negation} & && 
% \\ \midrule
% \textit{vʕalaʔ} & NPI `at all', used to reinforce negation & \sectref{tsova:reinforcing}& (\ref{tsova:flower2}–\ref{tsova:hair}) \\
% \textit{bedeⁿ} & NPI used in expletive negation expressing `only'& \sectref{tsova:nonneguse} & (\ref{tsova:octopus}) \\
% \textit{coħek'a} & negative conjunction `lest, otherwise' & \sectref{tsova:coordination} & (\ref{tsova:otherwise}) \\ 
% \lspbottomrule 
% \end{tabular}}
% \end{table}

\begin{table}
\caption{Summary of negation in Tsova-Tush}
\label{tsova:tab:sum}
\begin{tabularx}{\textwidth}{p{17mm}Qll}
\lsptoprule
\multicolumn{2}{l}{\textbf{Clausal negation}}                         & Chapter                           & Examples\\ 
\midrule
\textit{co} & preverbal particle used to negate     &                                   & \\
            & \hspace{4mm} (1) main clauses         & \ref{tsova:standardnegation}  & (\ref{tsova:horseshoe}–\ref{tsova:heardsong})\\
            & \hspace{4mm} (2) subordinate clauses  & \ref{tsova:nonmain}           & (\ref{tsova:sad}–\ref{tsova:care})\\
            & \hspace{4mm} (3) stative predications & \ref{tsova:stative}           & (\ref{tsova:ball}–\ref{tsova:donkey})\\
\textit{ma} & preverbal particle used to negate     &                           & \\
            & \hspace{4mm} (1) imperatives          & \ref{tsova:nondec}    & (\ref{tsova:pimple}--\ref{tsova:dontforget}) \\
            & \hspace{4mm} (2) optatives            & \ref{tsova:nondec}    & (\ref{tsova:seesun}) \\
    
\textit{cuj}& preverbal particle \textit{co} + interrogative enclitic \textit{=i} used to negate polar interrogatives & \ref{tsova:nondec} & (\ref{tsova:deserter}--\ref{tsova:CHILD}) \\

\textit{-gĕ} & verbal suffix only used in negative clauses, expresses `no longer' & \ref{tsova:standardnegation} & (\ref{tsova:anymore}–\ref{tsova:anymoreanymore})\\

\textit{ǯer} & particle to express `not yet' in negative clauses & \ref{tsova:standardnegation} & (\ref{tsova:yet}) \\

\textit{vuntxeʔ} & negative lexicalization `not know' & \ref{tsova:neglex} & (\ref{tsova:drankthat}) \\

\midrule 
\multicolumn{2}{l}{\textbf{Non-clausal negation}} &&\\
\midrule

\textit{co} & negative particle used as default negative reply to positive and negative questions & \ref{tsova:reply} & (\ref{tsova:no}--\ref{tsova:love})\\
\textit{moħec}& polarity reversing particle meaning `of course', used as affirmative reply to negative questions & \ref{tsova:reply} & (\ref{tsova:harmonica}) \\

\textit{co/ma}+ & {formation of negative pronouns} & {\ref{tsova:pronouns}} & {\tabref{tsova:tab:pron}} \\
interrogative pronoun &&&\\
\textit{-c'iⁿ} & suffix used for the formation of privatives from nouns and verbal stems & \ref{tsova:derivation} & (\ref{tsova:master}--\ref{tsova:uneducated}) \\
    
\textit{cuj} & negative particle \textit{co} + interrogative enclitic \textit{=i}, used to mark tag questions & \ref{tsova:nonneguse} & (\ref{tsova:poor}–\ref{tsova:respect}) \\

\midrule
\multicolumn{2}{l}{\textbf{Other aspects of negation}} && \\
\midrule

\textit{vʕalaʔ} & NPI `at all', used to reinforce negation & \ref{tsova:reinforcing}& (\ref{tsova:flower2}–\ref{tsova:hair}) \\
\textit{bedeⁿ} & NPI used in expletive negation expressing `only'& \ref{tsova:nonneguse} & (\ref{tsova:octopus}) \\
\textit{coħek'a} & negative conjunction `lest, otherwise' & \ref{tsova:coordination} & (\ref{tsova:otherwise}) \\
\lspbottomrule 
\end{tabularx}
\end{table}

\newpage%HC this is just so that the acknowledgments and references do not start before the summary table

\section*{Acknowledgements}

I'm thankful for the time and energy the Tsova-Tush speakers spend to teach me about their language, especially my main consultant Tina Tsiskarishvili.

\section*{Abbreviations}

\begin{tabularx}{.45\textwidth}{lQ}
\textsc{i-v} & gender\\
\textsc{1} & first person \\
\textsc{2} & second person \\
\textsc{3} & third person \\
\textsc{abl} & ablative \\
\textsc{abs} & absolutive \\
\textsc{adv} & adverbial \\
\textsc{all} & allative \\
\textsc{aor} & aorist\\
\textsc{caus} & causative\\
\textsc{cond} & conditional \\
\textsc{coord} & coordinating enclitic \\
\textsc{ccs} & contact case\\
\textsc{cvb} & converb\\
\textsc{dat} & dative\\
\textsc{dim} & diminutive\\
\textsc{dx} & deictic prefix\\
\textsc{erg} & ergative\\
\textsc{excl} & exclusive\\
\textsc{fut} & future\\
\textsc{gen} & genitive\\
\textsc{hab} & habitual\\
\textsc{imp} & imperative\\
\textsc{in} & spatial case `in'\\
\textsc{inf} & infinitive\\
\textsc{ins} & instrumental\\
\end{tabularx}
\begin{tabularx}{.45\textwidth}{lQ}
\textsc{ints} & intensifier\\
\textsc{intrz} & intransitivizer \\
\textsc{msd} & masdar \\
\textsc{neg} & negator\\
\textsc{obl} & oblique stem marker\\
\textsc{pl} & plural\\
\textsc{pol} & politeness \\
\textsc{pot} & potential \\
\textsc{priv} & privative\\
\textsc{proh} & prohibitive\\
\textsc{prs} & present\\
\textsc{pst} & past\\
\textsc{ptcp} & participle\\
\textsc{pvb} & preverb\\
\textsc{q} & polar question\\
\textsc{quot} & quotative \\
\textsc{refl} & reflexive\\
\textsc{rep} & reported past\\
\textsc{sg} & singular \\
\textsc{sim} & simultaneous\\
\textsc{sub} & subordinator\\
\textsc{subj} & subjunctive\\
\textsc{tr} & transitivizer\\
\textsc{voc} & vocative \\
\textsc{wit} & witnessed \\
\\
\end{tabularx}


\printbibliography[heading=subbibliography,notkeyword=this]

\end{document}
